\paragraph{Putting the residue extensions on one model.}
In the residue calculation write \(k'\) for the residue field and
use \(X\) for the characteristic-\(\ell\) cyclotomic parameter
throughout this calculation. Let \(D_i\) be the corresponding
étale \((\varphi,\Gamma)\)-module at letter \(i\), with Frobenius
\(\Phi_i\) coming from letter \(i-1\) and a principal cyclotomic
generator \(\gamma_i\) acting on letter \(i\).
Write \(\Phi_i N_{i-1}\) for the \(k'[[X]]\)-span of the
Frobenius images of a lattice. The convention for the residue
models after untwisting by \(\chi_1\) is
\[
 X^{w_i^*}N_i\subseteq \Phi_i N_{i-1}
       \subseteq X^{-w_i}N_i,\qquad
 (\gamma_i-1)N_i\subseteq XN_i,\qquad w_i^*=h_i-w_i.
\]
The model of the untwisted trivial quotient is the standard
rank-one lattice in its invariant basis. Its uniqueness among
the admissible models follows from the scalar comparison below.

\paragraph{Reduction of the marked Wach lattices.}
\label{par:residue-wach-realization}
Use the field, coefficient extension and reduced cyclotomic
realization of Section~\ref{subsec:cut-wach-models} with
\(p=\ell\), \(u=1+p\), and \(X\) the reduction of the integral
parameter \(X_0=-Z\). The residue field is \(k'\).
All crystalline and residual factors use the terminal Hodge
labels of that section.
The residual parameter actions, with Frobenius from the preceding letter, are
\[
 \varphi(X)=X^p,\qquad \gamma(X)\equiv X-X^2\pmod{X^3}.
\]


The first scalar character \(\chi_1\) has jump \(w_i\);
the scalar characters are \(\chi_c=\chi_1\epsilon^{1-c}\)
for \(1\le c\le t\), with \(\epsilon\) the cyclotomic character.
A full representation is crystalline with the labeled jumps
of the marked rank block together with those of the scalar
characters. Its residual characteristic polynomials agree with
those of that direct sum on the Galois group.

For a stable lattice
in such a full sum, take the Wach lattice after twisting by
\(\chi_1^{-1}\). With \(U_i\) the multiset of rank-block jumps before the twist, the full twisted multiset is
\[
 B_i=(U_i-w_i)\sqcup\{0,\ldots,t-1\}.
\]
Both parts retain their multiplicities. The separate rank and line lattices
are treated by the same twist. Their labeled height bounds,
after reduction by the coefficient uniformizer, give
\[
      X^{w_i^*}N_i\ \subset\ \Phi_i N_{i-1}
                    \ \subset\ X^{-w_i}N_i,
              \qquad (\gamma_i-1)N_i\subset XN_i .
\]
Here \(N_i\) is the reduction of the tame-invariant Wach lattice
inside the residual 'etale module on the corresponding letter.
A scalar line of sole jump \(e\) has Frobenius ideal \(X^e\).
For the trivial quotient line the invariant basis is \(1\).
\cite{berger-wach}\label{cor:residue-wach-realization-1}

\paragraph{The polynomial on a high-end letter.}
\begin{lemma}[Simple zero factor on a high-end letter]
\label{lem:residue-high-end-polynomial}
For the upper-end letter and \(s=w_i^*+1\), write \(D\) for the divided action on \(N_i/X^sN_i\). For \(P\) as in the proof the minimal polynomial satisfies
\[
 m_D(T)\mid TP(T),\qquad P(0)\ne0,\qquad \ker(D^j)=\ker(D)\quad(j\ge1).
\] Modulo \(X\)
the characteristic polynomial still records every jump residue
with its multiplicity.
\end{lemma}

\begin{proof}
In the integral tame parameter \(X_0\), with reduction \(X\),
put \(q=X_0+p\) and \(\nu=(\gamma-1)/X_0\). This is
\(O'\)-linear and preserves the invariant Wach lattice, since
\(X_0\) generates the invariant augmentation ideal and the
action is trivial on constants. For \(s\ge1\) define
\[
 I_s=\{a-b:\ 0\le a<s,\ b\in B_i\},\qquad
 d(m)=\frac{1-(1+p)^m}{p}\in\mathbb Z_p .
\]
On the characteristic-zero quotient by \(q^s\), the operator
is killed by
\[
 \prod_{m\in I_s}(T-d(m)),
\]
using each distinct
integer once. At the primitive level choose crystalline directions
adapted to the filtration. In a direction of jump \(b\) the
quotient has Laurent degrees \(-b,\ldots,s-1-b\) in the
local logarithmic parameter. The scalar substitution multiplies
this parameter by \(1+p\), and \(X_0=-p\) at the root.

In a direction of jump \(b\) the triangular operator \(T\) has distinct diagonal entries \(d(a-b)\), hence
\[
 \prod_{a=0}^{s-1}\big(T-d(a-b)\big)=0
 \quad\text{on this direction}.
\]
The product over the union of the eigenvalues thus kills every direction, including repeated jumps.
This is the finite-level filtration comparison of
Paragraph~\ref{par:wach-level-models}. The quotient of the
free invariant Wach lattice by \(q^s\) is free over \(O'\)
by monic division. The polynomial identity therefore holds
integrally. For \(s=1\) the characteristic polynomial has each
\(d(-b)\) with its multiplicity in \(B_i\), by the same
calculation and injectivity on coefficients.
\label{cor:residue-wach-realization-2}

Reduction gives \(q=X\), \(\bar d(m)=-m\), and operator
\[
 \nabla=(\gamma_i-1)/X.
\]
For negative \(m\) the congruence
for \(d(m)\) follows by inverting \((1+p)^{-m}\) modulo \(p^2\)
before division. The characteristic polynomial modulo \(X\)
therefore counts the residues of \(B_i\) with multiplicity.
At the upper end take \(s=w_i^*+1\). If \(0\le a<s\) and
\(a\equiv b\bmod p\), then \(a=b\). For a scalar or central
jump, writing \(b=u-w_i\), both \(u\) and \(w_i+a\) are in
the interval of width \(\eta'_i\) next to \(h_i\), including
the endpoint \(h_i\) for the latter.

For an outer jump the
congruence would give \(h_i-u\equiv w_i^*-a\bmod p\) in the
forbidden end window. It follows that
\[
   P(T)=\prod_{m\in I_s\setminus\{0\}}(T+m)\in k'[T],
   \qquad P(0)\ne0,\qquad
   \nabla P(\nabla)N_i\subset X^sN_i .
\]
If zero is absent from \(I_s\), \(P(\nabla)\) already has
image in \(X^sN_i\); applying \(\nabla\) gives the
displayed identity since the divided action preserves that
lattice. In the finite quotient, \(P(\nabla)\) is invertible
on the generalized zero space, since \(P(0)\ne0\). Thus
\(\nabla P(\nabla)=0\) there forces \(\nabla=0\). Using
the characteristic polynomial with all multiplicities would
lose this conclusion. \label{cor:residue-wach-realization-3}
\end{proof}

Finally an admissible residual representation is a quotient
of the reduction of its presenting full lattice. Its presenting
map is onto and kills \(\lambda T\), and untwisting preserves
that map. More generally take the intersection with a vector
subspace for a subrepresentation and the image for a quotient.
The intersection lattice is saturated. In a basis adapted to it,
Frobenius base change is exact, since the series ring is free
over its Frobenius image on \(1,X,\ldots,X^{p-1}\). Intersecting
the height inclusions or taking images gives both inclusions
for the induced lattices. The divided cyclotomic differences
intertwine the maps, so their polynomials restrict and descend:
for a stable vector subspace \(W\) one has
\[
 (X^sN_i)\cap W=X^s(N_i\cap W).
\]
The induced lattice sequence stays
exact modulo \(X\). The integral Wach intersection gives the
direct sum for a sum of lattices.

For the standard reference
use the split sum of the chosen rank lattice and the character
lines. Its reductions give the required models on those
summands, with their original labels. The prescribed residual
characteristic polynomials fix the composition factors of
each full sum, and untwisting multiplies them all by the
same character. This supplies the models used below with
their two heights, polynomials and target Sen counts.
\label{cor:residue-wach-realization-4}


\paragraph{A common lattice for the extension classes.}
\begin{lemma}[A largest compatible residue model]
\label{lem:largest-compatible-residue-model}
For two models \(N^0,N\) of the same residual representation,
each component is bounded by \(N_i\subseteq X^{-\sigma_i}N_i^0\). For the fixed residual representation the largest model is
\[
 N^{\max}=\sum_{N\ {\rm admissible}}N;
\]
the sum is itself admissible.
\end{lemma}

\begin{proof}
If \(N^0,N\) are two models of the same representation, let \(a_i\)
be the least nonnegative \(a\) with \(X^a N_i\subseteq N_i^0\).
The height inclusions give
\[
 \Phi_i N_{i-1}\subseteq X^{-a_i-w_i-w_i^*}\Phi_i N^0_{i-1}.
\]
In a Frobenius basis, the images of the coefficients have orders
multiplied by \(\ell\). Hence this inclusion implies
\(\ell a_{i-1}\le a_i+h_i\).
After \(m\) iterations this reads
\[
        a_i\le \ell^{-m}a_{i+m}
                 +\sum_{j=1}^m h_{i+j}/\ell^j .
\]
The first term tends to zero because the \(a_i\) form a
finite cyclic list. This gives \(a_i\le\sum_{j\ge1}h_{i+j}/\ell^j\).
Thus every model lies in \(X^{-\sigma_i}N_i^0\) in each letter.
The sum of two models is a model.

The sum model arises by direct sum
followed by the addition map of the common target. Exactness
and stability of the admissible subquotients give the induced
image lattice. The finite product of the containing lattices
is Noetherian over \(k'[[X]]\), so the sum of all models is
already a finite sum and is a largest model \(N^{\max}\).
This argument applies with the standard model as \(N^0\).
\label{cor:residue-model-length-1}
\end{proof}

In a residue extension, use the preimage lattice of the target
under the inclusion into the extension's model. Choose at each
letter a vector \(e_i\) mapping to \(1\) in the standard quotient.
The Herr pair on the target model is
\[
 b_i=\Phi_i e_{i-1}-e_i,\qquad g_i=(\gamma_i-1)e_i.
\]
Write \(P_i=\langle\Phi_i N_{i-1}\rangle\) for the target
Frobenius span. The cocycle equation and the two height tests
give
\[
 \begin{aligned}
 (\gamma_i-1)b_i&=\Phi_i g_{i-1}-g_i, &g_i&\in XN_i,\\
 X^{w_i}b_i&\in N_i &&\text{at a lower end},\\
 X^{w_i^*}b_i&\in P_i &&\text{at an upper end}.
 \end{aligned}
\]
Indeed \(\gamma_i\) acts trivially modulo \(X\) on the extension
lattice and \(N_i\) is the kernel there of the trivial quotient.
At the upper end apply the height inclusion to \(X^{w_i^*}e_i\).
Its image in the trivial line is \(X^{w_i^*}\), the coefficient
of \(\Phi_i e_{i-1}\) in the Frobenius span. Subtracting this
term leaves a vector in \(P_i\), by exactness on the induced
spans.

For the extra test at the upper end set \(s=w_i^*+1\).
Use the polynomial of Lemma~\ref{lem:residue-high-end-polynomial}
on the extension model, and put
\[
    e_i'=P(0)^{-1}P(\nabla)e_i,\qquad x_i=e_i-e_i'.
\]
The quotient map intertwines \(\nabla\) with zero, so \(e_i'\)
has the same quotient coordinate as \(e_i\), and \(x_i\in N_i\).
The identity \(\nabla P(\nabla)e_i\in X^s N_i^{\rm ext}\)
therefore yields
\[
     g_i-(\gamma_i-1)x_i=(\gamma_i-1)e_i'
                          \in X^{w_i^*+2}N_i .
\]
Here \(N_i^{\rm ext}\) denotes the extension lattice, and the
last containment uses its intersection with the target. Passing
to \(N^{\max}\) preserves the tested inclusions.
The exact functor on residue modules preserves addition of
extension classes and split extensions. A linear combination of
their Herr pairs that is a boundary on the common model thus
gives a coboundary for the original Galois classes.
\label{cor:residue-model-length-2}

\begin{auxlemma}[Sen multiplicities at the two ends]\label{lem:end-sen-multiplicities}
Use the residue positions and exclusions of Lemma~\ref{hyp:offwindow-local}.
Let \(s_i(y)\) be the generalized multiplicity of the residue Sen
operator at \(y\) on \(N_i^{\max}/XN_i^{\max}\) for the model of
Lemma~\ref{lem:largest-compatible-residue-model}.
Put \(d_i=\#\{u\in U_i:u<w_i\}\), with the labelled jumps counted
with multiplicity. At the indicated end letters one has
\[
 \begin{array}{ll}
 \displaystyle\sum_{y=-w_i}^{-1}s_i(y)=d_i &\text{at a lower end},\\[4pt]
 \displaystyle\sum_{y=0}^{w_i^*}s_i(y)=r-d_i &\text{at an upper end}.
 \end{array}
\]
\end{auxlemma}
\begin{proof}
Write \(N=N^{\max}\), and let \(N^0\) be the standard target
model. At letter \(i\), \(\nabla\) induces endomorphisms on
\(N_i/XN_i\) and \(N_i^0/XN_i^0\), with generalized
eigenvalue multiplicities \(s_i(y)\), \(s_i^0(y)\) at
\(y\in k'\). These are the Sen multiplicities in the
chosen normalization. For the standard target,
\(s_i^0(y)\) counts \(u\in U_i\) with
\[
 u-w_i\equiv y.
\]

The residue model comparison gives the same support condition
for the Sen weights of any target model, without a prescribed
multiplicity. Here is the support argument, which also works
for a subquotient of the marked rank block. The model induced on
that subquotient by \(N^0\) has Sen eigenvalues among \(u-w_i\)
for \(u\in U_i\). Indeed the subobject lattice is saturated and
the image lattice in the quotient is free. Exactness modulo \(X\)
and stability under \(\nabla\) give the weight inclusion.
For two models \(L,L'\) of that subquotient the carry bound
in both directions places \(L'/XL'\) as a subquotient of
\(X^{-\sigma_i}L/X^{\sigma_i+1}L\).

In particular each eigenvalue on \(L'/XL'\) lies in the set
\[
 \{b+a:b\text{ an eigenvalue on }L/XL,\ a\in\mathbb Z,\ |a|\le\sigma_i\}
\]
in the residue field.
On the other hand, any admissible model is induced from a
full crystalline sum, so its Sen weights are in the union of
the shifted rank-block jumps and the scalar residues
\(0,\ldots,t-1\). If an eigenvalue belongs to the scalar
block, the comparison with the induced model of \(N^0\)
would put a rank-block residue within \(\sigma_i\) of
that block. This is exactly the forbidden congruence from
the position tests. Hence all the model's Sen eigenvalues
come from the rank-block residues.
\label{cor:residue-model-length-3}

Let \(q_i(y)\) count generalized eigenvalues on \(N_i/N_i^0\).
On \(X^a N_i^0/X^{a+1}N_i^0\) the eigenvalues shift by \(-a\),
since the scalar action is normalized by
\[
 \gamma(X)=X-X^2\bmod X^3.
\]

In particular \(XN_i/XN_i^0\) has multiplicity \(q_i(y+1)\)
at \(y\). To compute the multiplicities in \(H=N_i/XN_i^0\), compare its
two filtrations. Their short exact sequences are
\[
 \begin{aligned}
 0&\longrightarrow XN_i/XN_i^0\longrightarrow H
        \longrightarrow N_i/XN_i\longrightarrow0,\\
 0&\longrightarrow N_i^0/XN_i^0\longrightarrow H
        \longrightarrow N_i/N_i^0\longrightarrow0 .
 \end{aligned}
\]
Adding generalized multiplicities along the two rows gives
\[
 s_i(y)+q_i(y+1)=s_i^0(y)+q_i(y).
\]
The filtration of \(X^{-\sigma_i}N_i^0/N_i^0\) shows that
if \(q_i(y)\ne0\), then \(y\equiv u-w_i+j\) for
some \(u\in U_i\) and \(1\le j\le\sigma_i\).
This step counts generalized eigenspaces, which have additive
dimensions in exact sequences even for nonsemisimple operators.

Let \(d_i\) count \(u\in U_i\) with \(u<w_i\).
The position tests give \(s_i(y)=0\) for
\[
 0\le y\le\sigma_i+1,
\]
and \(q_i(0)=0\).
Indeed a residue in these ranges would require \(u\)
congruent to \(w_i+y\) or \(w_i-j\) with distance
at most \(\sigma_i+2\) from the scalar block after checking
the end interval. For an outer jump the corresponding
congruence is precisely forbidden by the end-window test.

For a lower end letter sum the displayed identity for
\(-w_i\le y<0\). At the lower endpoint \(q_i(-w_i)=0\),
since any residue would require \(u\equiv -j\) with
\(1\le j\le\sigma_i\). For a central jump \(1\le u\le\eta'_i\), the end-width bound of
Lemma~\ref{hyp:offwindow-local} gives
\[
  0<u+j\le\eta'_i+\sigma_i<\ell ,
\]
which excludes \(u+j\equiv0\pmod\ell\). For an outer jump it contradicts
the residue exclusion. In the interval \([-w_i,0)\), the
standard multiplicities count exactly the central jumps below
\(w_i\). An outer jump cannot contribute, since its residue
would lie in the excluded lower end window. The endpoint vanishings \(q_i(0)=q_i(-w_i)=0\)
eliminate the correction in the telescoping identity:
\[
 \sum_{y=-w_i}^{-1}s_i(y)=
   \sum_{y=-w_i}^{-1}s_i^0(y)+q_i(-w_i)-q_i(0)=d_i.
\]

For an upper end letter sum from \(0\) through
\(w_i^*+\sigma_i+1\). At the upper endpoint after telescoping,
\(q_i(w_i^*+\sigma_i+2)=0\): a contributing \(u\) would
be congruent to \(h_i+\sigma_i+2-j\), \(1\le j\le\sigma_i\).
For an outer jump this is precisely the residue exclusion.

At this end a central jump satisfies \(h_i-\eta'_i\le u\le h_i-1\), so
\[
 0<h_i+\sigma_i+2-j-u\le\eta'_i+\sigma_i+1<\ell .
\]
This excludes the congruence also for a central jump. The same test gives
\(s_i(y)=0\) for \(w_i^*<y\le w_i^*+\sigma_i+1\).
Throughout the summation range the standard multiplicities
count exactly the jumps \(u\ge w_i\), all of which are in
the upper end interval. Put \(B=w_i^*+\sigma_i+1\). The vanishing of
\(s_i(y)\) for \(w_i^*<y\le B\) gives
\[
 \sum_{y=0}^{w_i^*}s_i(y)=\sum_{y=0}^{B}s_i(y)
     =\sum_{y=0}^{B}s_i^0(y)+q_i(0)-q_i(B+1)=r-d_i .
\]
All the interval counts include multiplicities and use bounds
strictly shorter than \(\ell\) at the relevant end. The number
\(d_f\) in the local comparison is \(\sum_i d_i\).
\label{cor:residue-model-length-4}
\end{proof}

\paragraph{From the polar parts to the residue rank.}
The filtration of \(N_i/X^{\sigma_i+2}N_i\) has successive
eigenvalues shifted by \(0,-1,\ldots,-(\sigma_i+1)\).
The Sen vanishing in the proof of Lemma~\ref{lem:end-sen-multiplicities}
makes \(\nabla\) invertible on this quotient.
For each pair choose \(y_i\in N_i\) so that
\[
 g'_i=g_i-(\gamma_i-1)y_i\in X^{\sigma_i+3}N_i,\qquad
 b'_i=b_i-\Phi_i y_{i-1}+y_i .
\]
This is a change by a Herr boundary.
Set \(c_i=b_i-\Phi_i y_{i-1}+x_i\), using \(x_i=y_i\)
at a lower end letter and the deeper cyclotomic correction
at an upper end letter.

The corrected vector \(c_i\) has the same image as \(b'_i\)
modulo \(N_i\).
The relation \(\ell(\sigma_{i-1}+1)\ge h_i+\sigma_i+1\)
follows by shifting the carry series. Consequently
\(\Phi_i g'_{i-1}\in X^{w_i^*+2}N_i\).
The height inclusions and the cocycle equation give
\[
\begin{array}{ll}
c_i\in X^{-w_i}N_i,&(\gamma_i-1)c_i\in XN_i
       \quad\text{at a lower end letter},\\
c_i\in X^{-w_i^*}\Phi_i N_{i-1},&
(\gamma_i-1)c_i\in X^{w_i^*+2}N_i
       \quad\text{at an upper end letter}.
\end{array}
\]
If a linear combination of the \(c_i\) is integral at
every letter, the combination of the \(b'_i\) is integral.
Its cyclotomic difference belongs to \(X^{\sigma_i+3}N_i\),
using the deeper \(g'_i\) and the same carry inequality.
Invertibility modulo \(X^{\sigma_i+2}\) places \(b'_i\) in
\(X^{\sigma_i+2}N_i\).

On the product of these deep lattices
\(\Phi\) strictly increases the order at each step by the
height and carry inequalities. Solve \((\Phi-1)v=b'\) there
by \(-\sum_{a\ge0}\Phi^a b'\). Commutation with \(\gamma\)
and injectivity of \(\Phi-1\) on the deep product give
\[
 (\gamma-1)v=g'.
\]
Thus the original combination is a
Herr boundary. In particular the dimension of a space of
independent pairs is at most the sum of the bounds for the
images of their \(c_i\)'s modulo \(N_i\).

At a lower end letter the possible images lie in the kernel
of \(\nabla\) on \(X^{-w_i}N_i/N_i\).
On its graded piece of degree \(a\) the generalized zero
multiplicity is \(s_i(a)\). The sum for \(-w_i\le a<0\)
is \(d_i\), so the kernel has dimension at most \(d_i\).

At an upper end letter put \(M=\Phi_i N_{i-1}\),
\[
 V=X^{-w_i^*}M
\]
and \(J=X^{w_i^*+1}N_i\).
Then \(J\subseteq N_i\subseteq V\), and \(V/J\) is finite
dimensional over \(k'\). The required image is that of
\[
 \ker(\nabla:V/J\longrightarrow D_i/J)
\]
in \(D_i/N_i\).
The codimension estimate for \(\nabla\) on \(V/J\)
uses the cyclotomic Frobenius left inverse \(\psi_i\).
In the range \(1\le w_i^*\le\ell-7\), its scalar companion
maps \(X^{1-w_i^*}k'[[X]]\) into \(k'[[X]]\), and
\(\psi_i\Phi_i=1\). On the part of \(X^{1-w_i^*}M\) killed
by \(\psi_i\), the operator \(\gamma_i-1\) can be inverted
in \(X^{-w_i^*}M\).

\begin{lemma}[Frobenius coefficient and divided-action inverse]
\label{lem:residue-frobenius-left-inverse}
For an upper-end letter, set \(d=w_i^*\) with \(1\le d\le\ell-7\)
and \(M=\langle\Phi_i N_{i-1}\rangle\) in the residue model of
Lemma~\ref{lem:largest-compatible-residue-model}. Write \(\Psi_i\) for the Frobenius left inverse and \(\Psi\) for its scalar counterpart. Then
\[
 \Psi_i(X^{1-d}M)\subseteq N_{i-1},\qquad\Psi(X^{1-d}k'[[X]])\subseteq k'[[X]].
\]
On the kernel in \(X^{1-d}M\) the \(\gamma_i-1\) equation
has a solution in \(X^{-d}M\) in the kernel.
\end{lemma}

\begin{proof}
Extend scalars temporarily to
\(k'((\pi))\), with \(X\) the invariant parameter of
Paragraph~\ref{par:residue-wach-realization}.
Write \(p=\ell,\ d=w_i^*\) and \(A=1+\pi\).
Use the extended lattices \(\widetilde N=k'[[\pi]]N_{i-1}\)
and \(\widetilde M=k'[[\pi]]M\). A basis of \(\widetilde N\)
from \(N_{i-1}\) gives a Frobenius basis of \(\widetilde M\).

In this basis every vector in the ambient space has a unique
expansion
\[
             \sum_{j=0}^{p-1} A^j\Phi_i(y_j).
\]
Indeed \(1,\pi,\ldots,\pi^{p-1}\) form the usual basis over
the image under \(\pi\mapsto\pi^p\). The triangular change
to \(1,A,\ldots,A^{p-1}\) is integral with unit diagonal.
Define \(\psi_i\) to return \(y_0\). This definition is
independent of the Frobenius basis, by uniqueness of the
direct sum expansion. A tame substitution of exponent \(z\)
sends the \(j\)-summand to the summand with index \(zj\bmod p\),
since
\[
 A^{zj}=A^{j'}\phi(A^{(zj-j')/p}).
\]
Thus the
projection descends to the tame invariants, commutes with
the principal action and satisfies \(\psi_i\Phi_i=1\).
The same expansion in rank one gives the scalar companion.

If the input lies in \(X^{1-d}M\), its ordinary order in
\(\widetilde M\) is at least \((p-1)(1-d)\ge p(1-d)\).
The integral change of basis gives
\(y_j\in\pi^{1-d}\widetilde N\) for every \(j\).
The zeroth component is tame-invariant. Its nonzero
coefficients have orders divisible by \(p-1\) in the
invariant basis. Since \(1-d>-(p-1)\), it belongs to
\(N_{i-1}\). This proves in particular the scalar bound.

For an input with \(y_0=0\), solve separately for \(1\le j<p\).
On the \(j\)-summand the difference is
\[
   (\gamma_i-1)\bigl(A^j\Phi_i(z)\bigr)
      = A^j\Phi_i\bigl((A^j\gamma_{i-1}-1)z\bigr),
\]
since \(\gamma(A)=A^{1+p}\). The operator on the right is
multiplication by \(A^j-1\), of ordinary order one, plus
\(A^j(\gamma_{i-1}-1)\). On each
\(\pi^h\widetilde N\), \(h\in\mathbb Z\), the latter term
increases order by at least \(p-1\): use
\(\gamma(\pi)/\pi\equiv1\bmod\pi^{p-1}\) and
\((\gamma_{i-1}-1)N_{i-1}\subset XN_{i-1}\).
After dividing by \(A^j-1\), the error raises order by at
least \(p-2>0\). Put \(B_j=(A^j-1)^{-1}A^j(\gamma_{i-1}-1)\).

The operator \(B_j\) increases order by at least \(p-2\), so the solution is
\[
  z_j=\sum_{k\ge0}(-B_j)^k(A^j-1)^{-1}y_j
                       \ \in\ \pi^{-d}\widetilde N .
\]
Multiplication by \(A^j-1\) after applying \(1+B_j\)
gives \((A^j\gamma_{i-1}-1)z_j=y_j\).
The resulting sum has ordinary order at least \(-pd\) in
\(\widetilde M\). Tame averaging preserves that order, the
equation and the vanishing under \(\psi_i\). For the averaged
vector, the lower order bound rounds up to \(-(p-1)d\), since
\(d<p-1\) and the coefficients in the invariant basis have
orders in \((p-1)\mathbb Z\). This gives the preimage in
\(X^{-d}M\) with the required one-power loss.
\label{cor:residue-frobenius-inverse-1}
\end{proof}



By Lemma~\ref{lem:residue-frobenius-left-inverse}, for \(u\in V\) apply \(\psi_i\) to \(Xu\), obtaining
\(y\in N_{i-1}\). Subtract \(\Phi_i y\) from \(Xu\)
and use the inverse on the \(\psi_i=0\) part.
Divide by \(X\) to express \(u\) as
\(X^{-1}\Phi_i y+\nabla v\), with \(v\in V\).
In a basis of \(N_{i-1}\), write \(y=y_0+Xy_1\)
with \(y_0\) in the span of the constant basis vectors.
The preceding letter's Sen vanishing solves
\[
 \nabla z\equiv y_1\bmod X^{\sigma_{i-1}+2}N_{i-1}.
\]

Then \(X^{-1}\Phi_i(Xy_1)\) differs from
\(\nabla\Phi_i z\) by an element of \(J\), by commutation
and the carry inequality. It follows that \(V\) is covered
by \(\nabla V+J\) and the span of the \(r\) vectors
\(X^{-1}\Phi_i\) of the preceding basis.
Rank-nullity gives dimension at most \(r\) to
\(\ker(\nabla:V/J\to D_i/J)\), since \(V/J\) lies in the
sum of that \(r\)-dimensional contribution and the image
of this map.

The zero generalized eigenspace on \(N_i/J\) has dimension
\[
 \sum_{a=0}^{w_i^*}s_i(a)=r-d_i
\]
by its graded pieces.
The polynomial with \(P(0)\ne0\) gives
\(\nabla P(\nabla)=0\) on \(N_i/J\). Hence on the generalized
zero eigenspace \(\nabla=0\), since \(P(\nabla)\) is invertible
there. These \(r-d_i\) directions of the kernel have zero
image modulo \(N_i\). The image has dimension at most \(d_i\).
This proves the rank bound \(\sum_i d_i=d_f\) on the
common residue model.
\label{cor:residue-model-length-5}

\paragraph{The residue invariants and the scalar lines.}
\begin{proposition}[The scalar lines on the residue model]
\label{prop:offwindow-scalar-residue-model}
In the residue model of Lemma~\ref{lem:largest-compatible-residue-model},
with heights and exclusions of Lemma~\ref{hyp:offwindow-local}, none of \(C_e=\bar\epsilon^{-e}\) for
\(-1\le e<t\) occurs as a subquotient of the marked rank block.
The marked rank block therefore has zero local cohomology in
degrees zero and two. For \(0\le e<t\), the induced lattice on a scalar
subquotient \(C_e\) is its standard model. If \(A_e\) is the space of admissible residue extension classes with trivial quotient and target \(C_e\), then
\[
 \dim A_0\le1,\qquad A_e=0\quad(1\le e<t).
\]
\end{proposition}

\begin{proof}
\proofstep{The carry estimate and the excluded scalar residues.}
Let \(\bar\epsilon\) be the residue cyclotomic character and
write \(C_e\) for the rank-one representation
\(\bar\epsilon^{-e}\). The standard model \(L_i^{(e)}\) has
Sen eigenvalue \(e\) and its Frobenius span is \(X^eL_i^{(e)}\).
For \(-1\le e<t\) these lattices satisfy the two height
inclusions. Throughout this range the endpoints satisfy
\[
 w_i\ge1,\qquad w_i+e\le h_i-1.
\]
The standard models can therefore
serve as comparison lattices whether or not the line
is a subquotient under study.

Suppose \(C_e\) were a subquotient of the marked rank block
for such an \(e\). The induced lattice on that subquotient
has a Sen eigenvalue congruent to \(e+a\) with
\(|a|\le\sigma_i\), by the two-direction carry comparison
with the standard line. Its eigenvalue comes from the
Sen weights of \(N_i^0\) by exactness.
The residue exclusion near the scalar block forbids
\[
 u\equiv w_i+e+a
\]
for \(u\in U_i\), since this point
is within \(\sigma_i+2\) of the scalar interval.
For \(e=0\) this rules out invariants. For \(e=-1\)
it rules out a quotient of cyclotomic character,
and local Tate duality gives vanishing in degree two.
\label{cor:residue-model-length-6}

\proofstep{Uniqueness of the induced scalar lattice.}
For \(0\le e<t\) the standard line is the unique
model of \(C_e\) induced from the full crystalline sums.
To see this, refine the subquotient presentation of
the line to a composition series of the full residue
representation. Use the induced lattices at each step.
At the level of semisimplifications, one full copy has the same factors as
\[
 \bar W^{\rm ss}\ \oplus\ \bigoplus_{e=0}^{t-1} C_e,\qquad C_e=\bar\epsilon^{-e},
\]
where \(\bar W\) denotes the marked residual rank block after twisting by \(\bar\chi_1^{-1}\). Direct sums repeat this list with multiplicity.
The preceding exclusion keeps the two lists disjoint, and
the scalar characters are distinct since \(t<\ell-1\).
The Sen weights of a factor from the rank block are among
its shifted jump residues by the support argument above.

A lattice on a scalar factor \(C_a\) is \(X^{\kappa_i}L_i^{(a)}\)
for an integer \(\kappa_i\). The two-direction bound gives
\[
 |\kappa_i|\le\sigma_i.
\]
Its height conditions read
\[
 -w_i+\kappa_i\le \ell\kappa_{i-1}+a
                      \le w_i^*+\kappa_i .
\]
The end-width bound of Lemma~\ref{hyp:offwindow-local}
gives \(\eta'_i+\sigma_i<\ell\). Use the left height inequality
at a lower end, and the right one at an upper end, to obtain
\begin{equation}\label{eq:scalar-end-signs}
 \begin{array}{ll}
 \ell\kappa_{i-1}\ge -(w_i+a)+\kappa_i
          \ge-\eta'_i-\sigma_i>-\ell
     &\text{at a lower end},\\
 \ell\kappa_{i-1}\le (w_i^*-a)+\kappa_i
          \le\eta'_i+\sigma_i<\ell
     &\text{at an upper end}.
 \end{array}
\end{equation}

At a fixed letter \(i\), index the scalar composition factors by
\(a_\nu\), with repetitions. Every shift \(\kappa_{i,\nu}\)
is therefore nonnegative if the next letter is a lower end and
nonpositive if it is an upper end. All shifted weights
\(a_\nu-\kappa_{i,\nu}\) belong to
\([-\sigma_i,t-1+\sigma_i]\). This interval has length less than
\(\ell\), contains the nominal scalar weights and excludes the
rank-block weights. Additivity of generalized Sen multiplicities
therefore gives equality of integer multisets in this interval.
In particular
\[
 \sum_\nu (a_\nu-\kappa_{i,\nu})=\sum_\nu a_\nu
 \quad\Longrightarrow\quad\sum_\nu \kappa_{i,\nu}=0 .
\]
The sign from \eqref{eq:scalar-end-signs} forces each shift to vanish. This also applies to the line in
the original subquotient presentation.
\label{cor:residue-model-length-7}

\proofstep{Removing pole terms in an admissible Herr pair.}
Now take an admissible extension with target \(C_e\),
\(0\le e<t\), and trivial quotient. Its Herr pair \(b_i,g_i\)
is on these standard lattices. We have \(g_i\in XL_i^{(e)}\).
At a lower end letter \(b_i\in X^{-w_i}L_i^{(e)}\), and
at an upper end letter
\[
 b_i\in X^{e-w_i^*}L_i^{(e)}.
\]

In either case the permitted pole degrees \(j>0\)
satisfy \(e+j<\ell\), by the short end width:
\(w_i+e\le\eta'_i\) on the lower side and
\(w_i^*\le\eta'_i\) on the upper side.
The cocycle equation implies \(\nabla b_i\in L_i^{(e)}\),
since Frobenius takes \(XL_{i-1}^{(e)}\) to
\(X^{\ell+e}L_i^{(e)}\). On the negative degree \(-j\)
the leading eigenvalue of \(\nabla\) is \(e+j\ne0\).
Inductively eliminating the negative coefficients shows
\[
 b_i\in L_i^{(e)}.
\]
For \(e>0\), \(\Phi\) is contracting
on \(\prod_iL_i^{(e)}\). Inverting \(\Phi-1\) there kills
the pair by the same cochain argument as above.

\proofstep{The constant class on the cycle.}
For \(e=0\), reduce \(b_i\) modulo \(X\) to a constant
\(c_i\) in the invariant basis. Invert \(\Phi-1\)
on \(\prod_i XL_i^{(0)}\) for \(b_i-c_i\).
The cocycle equation again identifies the other
coordinate as the cyclotomic difference of the resulting
vector, since \(\gamma_i\) fixes the \(c_i\).
Let \(q\) be the number of letters and \(k\) the residue
field. The last coordinate space is a subquotient of
\[
 k^q/\{(d_{i-1}-d_i)_i : d\in k^q\}.
\]
On the cycle every difference vector has zero sum. Conversely
if \(\sum_i c_i=0\), start with \(d_0=0\) and solve
\(d_i=d_{i-1}-c_i\) successively. The last equation closes
because the sum is zero. Thus summation identifies the displayed
quotient with \(k\). This calculation uses no division by \(q\)
and gives a dimension bound of one in every residue characteristic.
\label{cor:residue-model-length-8}
\end{proof}
